\loai{Tính $N_2$}
\begin{vidu} %[3P3-14.3K][Loai2] 
	Đặt vào hai đầu cuộn sơ cấp của máy biến áp lí tưởng điện áp xoay chiều có giá trị hiệu dụng không đổi. Nếu quấn thêm vào cuộn thứ cấp 90 vòng thì điện áp hiệu dụng hai đầu cuộn thứ cấp để hở thay đổi $30\%$ so với lúc đầu. Số vòng dây ban đầu ở cuộn thứ cấp là
	\choice
	{1200 vòng}
	{\True 300 vòng}
	{900 vòng}
	{600 vòng}
	\loigiai{
		\Binhluan{
		\begin{align*}
			&\dfrac{N_2}{N_1}=\dfrac{U_2}{U_1}\\ 
			&\dfrac{{N_2}+90}{N_1}=\dfrac{{U_2}+0,3{U_2}}{U_1}=1,3. \dfrac{U_2}{U_1}\\
			\Rightarrow\ &\dfrac{{N_2}+90}{N_2}=1,3\Rightarrow \dfrac{90}{N_2}=0,3\Rightarrow {N_2}=300\ \ct{vòng}.
		\end{align*}
	}
{
Quấn thêm vào cuộn thứ cấp thì $N_2$ tăng $\Rightarrow$ $U_2$ tăng.
}
	}
\end{vidu}
\loai{Thay đổi $N_2-$ Tính $U_2$}
\begin{vidu} %[3P3-14.3B][Loai3] 
	Đặt vào hai đầu cuộn sơ cấp của máy biến áp lí tưởng một điện áp xoay chiều có giá trị hiệu dụng không đổi thì điện áp hiệu dụng hai đầu cuộn thứ cấp để hở là 400 V. Nếu giảm bớt số vòng dây của cuộn thứ cấp đi một nửa so với ban đầu thì điện áp hiệu dụng hai đầu cuộn thứ cấp là
	\choice
	{100 V}
	{\True 200 V}
	{600 V}
	{800 V}
	\loigiai{
	\Binhluan{
		Áp dụng công thức của máy biến áp, ta có hệ
		\begin{align*}
		 \begin{cases}
		\dfrac{U_1}{400}=\dfrac{N_1}{N_2}\\
		\dfrac{U_1}{U'_2}=\dfrac{N_1}{0,5N_2}\\
		\end{cases}\Rightarrow U_2'=0,5.400=200\ V.
	\end{align*}
}
{Mấu chốt là $U_1=const.$}	
}
\end{vidu}
\begin{vidu} %[3P3-14.3K][Loai3] 
	Đặt vào hai đầu cuộn sơ cấp của máy biến áp lí tưởng một điện áp xoay chiều có giá trị hiệu dụng không đổi thì điện áp hiệu dụng hai đầu cuộn thứ cấp để hở là 300 V. Nếu giảm bớt một phần ba số vòng dây của cuộn thứ cấp thì điện áp hiệu dụng hai đầu của nó là
	\choice
	{100 V}
	{220 V}
	{\True 200 V}
	{110 V}
	\loigiai{
		\Binhluan{
		Áp dụng công thức máy biến áp cho hai trường hợp:
		\begin{align*}
			\begin{cases}
				\dfrac{300}{U_1}=\dfrac{N_2}{N_1}\\
				\dfrac{U_2^{'}}{U_1}=\dfrac{N_2-\dfrac{N_2}{3}}{N_1}
			\end{cases}\Rightarrow \begin{cases}
				\dfrac{300}{U_1}=\dfrac{N_2}{N_1}\\
				\dfrac{U_2^{'}}{U_1}=\dfrac{2}{3}.\dfrac{N_2}{N_1}
			\end{cases}\Rightarrow U_2^{'}=\dfrac{2}{3}.300=200\ V.
		\end{align*}
	}
{Để cho gọn ta có thể đặt $k=\dfrac{N_2}{N_1}$ và chuẩn hoá $U_1=1$}
}
\end{vidu}
\loai{Thay đổi $N_2-$ Tính $U_1$}
\begin{vidu} %[3P3-14.3K][Loai4] 
	Đặt vào hai đầu cuộn sơ cấp của máy biến áp lí tưởng một điện áp xoay chiều có giá trị hiệu dụng không đổi. Cuộn sơ cấp có 200 vòng dây. Khi máy biến áp hoạt động người ta đo được điện áp hiệu dụng trên hai đầu dây của cuộn thứ cấp là 100 V. Nếu quấn thêm vào cuộn thứ cấp thêm 10 vòng dây thì điện áp hiệu dụng đo được trên cuộn thứ cấp là 120 V. Điện áp hiệu dụng trên hai đầu cuộn sơ cấp là
	\choice
	{200 V}
	{\True 400 V}
	{250 V}
	{300 V}
	\loigiai{
\Binhluan{
		\begin{itemize}
			\item Ta có: $\dfrac{U_2}{U_1}=\dfrac{N_2}{N_1}$.
			\item Theo bài ra: $\left\{{
				\begin{array}{l}
				{\dfrac{100}{U_1}=\dfrac{N_2}{200}}\\
				{\dfrac{120}{U_1}=\dfrac{N_2+10}{200}}\\
				\end{array}}\right.\Rightarrow \dfrac{100}{120}=\dfrac{N_2}{N_2+10}\Rightarrow N_2=50$ vòng
			$\Rightarrow U_1=\dfrac{U_2.N_1}{N_2}=\dfrac{100.200}{50}=400\ V$.
	\end{itemize}
}
{
Đây là bài toán tường minh với hệ 2 phương trình 2 ẩn nên không thể chuẩn hoá đại lượng nào.
}
}
\end{vidu}
\loai{Thay đổi cả $N_1$ và $N_2$}
\begin{vidu} %[3P3-14.3K][Loai5] 
	Cuộn sơ cấp của một máy biến áp lí tưởng có $N_1$ vòng dây. Khi đặt một điện áp xoay chiều có giá trị hiệu dụng 120 V vào hai đầu cuộn sơ cấp thì điện áp hiệu dụng ở hai đầu cuộn thứ cấp để hở đo được là 100 V. Nếu tăng thêm 150 vòng dây cho cuộn sơ cấp và giảm 150 vòng dây ở cuộn thứ cấp thì khi đặt vào hai đầu cuộn sơ cấp một điện áp hiệu dụng 160 V thì điện áp hiệu dụng ở hai đầu cuộn thứ cấp để hở vẫn là 100 V. Số vòng dây ở cuộn sơ cấp ban đầu $N_1$ là
	\choice
	{$N_1=825$ vòng}
	{$N_1=1320$ vòng}
	{\True $N_1=1170$ vòng}
	{$N_1=975$ vòng}
	\loigiai{
		\Binhluan{
		Theo giả thiết bài toán, ta có:
		\begin{align*}
			&\begin{cases}
				\dfrac{N_1}{N_2}=\dfrac{120}{100}\\
				\dfrac{N_1+150}{N_2-150}=\dfrac{160}{100}
			\end{cases}\Rightarrow \begin{cases}
				N_2=\dfrac{5}{6}.N_1\\
				\dfrac{N_1+150}{N_2-150}=\dfrac{8}{5}
			\end{cases}\\
			\Rightarrow\  &\dfrac{N_1+150}{\dfrac{5}{6}.N_1-150}=\dfrac{8}{5}\xrightarrow{Shift\to Solve}N_1=1170.
		\end{align*}
	}
{Ở bài toán này ta lại thay đổi $U_1$ và giữ nguyên $U_2$}
	}
\end{vidu}
\loai{Thay đổi $N_2$ qua 2 giá trị - Tính $U_2$}
\begin{vidu} %[3P3-14.3K][Loai7] 
	Đặt vào hai đầu cuộn sơ cấp của một máy biến áp lí tưởng một điện áp xoay chiều có giá trị hiệu dụng không đổi thì hiệu điện thế hiệu dụng giữa hai đầu mạch thứ cấp khi để hở là 200 V. Khi ta giảm bớt n vòng dây ở cuộn sơ cấp thì hiệu điện thế hiệu dụng giữa hai đầu mạch thứ cấp khi để hở là U; nếu tăng n vòng dây ở cuộn sơ cấp thì hiệu điện thế hiệu dụng giữa hai đầu mạch thứ cấp khi để hở là 0,5U. Giá trị của U là
	\choice
	{250 V}
	{200 V}
	{100 V}
	{\True 300 V}
	\loigiai{
		\Binhluan{
		\begin{itemize}
			\item Ban đầu : $\dfrac{N_1}{N_2}=\dfrac{U_1}{200}$.
			\item Khi ta giảm bớt n vòng dây ở cuộn sơ cấp:  $\dfrac{{N_1}-n}{N_2}=\dfrac{U_1}{U}$.
			\item Khi tăng n vòng dây ở cuộn sơ cấp 
			\begin{align*}
				&\dfrac{{N_1}+n}{N_2}=\dfrac{2{U_1}}{U}
				\Rightarrow 2\left( \dfrac{{N_1}-n}{N_2} \right)=\dfrac{{N_1}+n}{N_2}\Leftrightarrow 2\left( {N_1}-n \right)={N_1}+n\\
				\Leftrightarrow &{N_1}=3n
				\Rightarrow\ \dfrac{{N_1}-n}{N_2}=\dfrac{U_1}{U}\Leftrightarrow \dfrac{2{N_1}}{3{N_2}}=\dfrac{U_1}{U}\Leftrightarrow \dfrac{2}{3}\left( \dfrac{U_1}{200} \right)=\dfrac{U_1}{U}\Leftrightarrow U=300\ V.
			\end{align*}
		\end{itemize}
	}
{
Để giải hệ phương trình khéo léo ta có thể đặt $k=\dfrac{N_2}{N_1}$ và chuẩn hoá $U_1=1$.
}	
	}
\end{vidu}
\loai{Tính số vòng dây quấn thêm}
\begin{vidu} %[3P3-14.3K][Loai9] 
	\nguon{ĐH - 2011} Một học sinh quấn một máy biến áp với dự định số vòng dây của cuộn sơ cấp gấp hai lần số vòng dây của cuộn thứ cấp. Do sơ suất nên cuộn thứ cấp bị thiếu một số vòng dây. Muốn xác định số vòng dây thiếu để quấn tiếp thêm vào cuộn thứ cấp cho đủ, học sinh này đặt vào hai đầu cuộn sơ cấp một điện áp xoay chiều có giá trị hiệu dụng không đổi, rồi dùng vôn kế xác định tỉ số điện áp ở cuộn thứ cấp để hở và cuộn sơ cấp. Lúc đầu tỉ số điện áp bằng 0,43. Sau khi quấn thêm vào cuộn thứ cấp 24 vòng dây thì tỉ số điện áp bằng 0,45. Bỏ qua mọi hao phí trong máy biến áp. Để được máy biến áp đúng như dự định, học sinh này phải tiếp tục quấn thêm vào cuộn thứ cấp
	\choice
	{\True 60 vòng dây}
	{84 vòng dây}
	{100 vòng dây}
	{40 vòng dây}
	\loigiai{
		\Binhluan{
		\begin{itemize}
			\item Gọi $N_1$, $N_2$ là số vòng dây ban đầu của mỗi cuộn; n là số vòng phải cuốn thêm cần tìm.
			\item Ta có: $\dfrac{{N}_{2}}{{N}_{1}}=0,43;\dfrac{{N}_{2}+24}{{N}_{1}}=0,45\Rightarrow 
			\begin{cases}
			N_{1}=1200\\
			N_{2}=516
			\end{cases}$
			\item Do $\dfrac{{N}_{1}}{{N}_{2}+24+n}=2\Rightarrow n=60$.
	\end{itemize}
}
{Đề bài hỏi số vòng quấn thêm. Cần để ý kĩ tiến trình thí nghiệm và lập nhanh các phương trình tương ứng với dữ kiện đề bài}
}
\end{vidu}
\loai{Quấn ngược}
\begin{vidu} %[3P3-14.3K][Loai10]
	Một người định cuốn máy biến áp có điện áp hiệu dụng ngõ vào (cuộn sơ cấp) là $U_1=220$ V và điện áp hiệu dụng muốn đạt được ở ngõ ra (cuộn thứ cấp) là $U_2=24$ V. Xem máy biến áp là lí tưởng. Các tính toán về mặt kĩ thuật cho kết quả cần phải quấn 1,5 (vòng/vôn). Người đó cuốn đúng hoàn toàn cuộn sơ cấp nhưng lại cuốn ngược chiều những vòng cuối của cuộn thứ cấp. Khi thử máy với điện áp sơ cấp là 110 V thì điện áp thứ cấp đo được 10 V. Số vòng dây bị cuốn ngược chiều là
	\choice
	{12}
	{20}
	{\True 3}
	{6}
	\loigiai{
	\Binhluan{
		\begin{itemize}
			\item Nếu quấn đúng thì số vòng sơ cấp và thứ cấp của máy biến áp là $\begin{cases}
			N_1=1,5.220=330\\
			N_2=1,5.24=36\\
			\end{cases}$
			\item Giả sử người đó quấn ngược n vòng thì từ trường cảm ứng trong n vòng này sẽ triệt tiêu từ trường cảm ứng của n vòng quấn đúng, do vậy $\dfrac{10}{110}=\dfrac{36-2n}{330}\Rightarrow n=3$.
		\end{itemize}
	}
{
$q=1,5$ vòng/vôn: $N=Uq$.\\
Bài toán \textbf{quấn ngược} đòi hỏi ta phải tư duy nhanh; số vòng quấn ngược vừa mất thêm n vòng cuốn mà lại huỷ bớt n vòng nên số vòng của nó tham gia vào công thức là $N-2n$.
}
	}
\end{vidu}
